% Validation subsection (2.6) and Section 3, \input by main.tex.
% Every number must trace to the CSVs the pipeline notebooks write
% (tide_boundary_comparison_<site>.ipynb sections 10/10b,
% intertidal_topography_villaviciosa.ipynb sections 9a/9b); the scripts
% experiments/v4_lag_external.py, v5_rtk_onsite.py and v6_lidar_external.py
% only assemble them into results/*.json and the paper figures.
% Convention: "MAREA's lag at the inner gauge" is the smoothed band profile
% evaluated at the gauge's distance from the mouth, i.e. what a pixel there
% receives; the band that contains the gauge and its 95 % interval are
% given beside it. The comparison method's lag is not scored against gauges.
% Runs of 2026-09-23 (fixed-band product): Ferrol, Westerschelde, Ems and
% Wadden done; Villaviciosa (RTK, LiDAR) pending -> \ldots.

\subsection{Validation}\label{sec:validation}

The method yields two quantities, the lag of each band and the elevation
of each pixel, each validated on site, in the R\'ia de Villaviciosa, and
externally, against tide-gauge records and surveys made by others
(Table~\ref{tab:validation}). The lag blocks ask whether the measured
lag is real and whether applying it improves the level assigned inside
the estuary; the elevation blocks set both maps against independent
surveys. Nothing that judges a result enters a fit.

\begin{table*}[t]
\centering
\caption{The four validation blocks: what is validated, against what,
where, and with which metric.}
\label{tab:validation}
\footnotesize
\setlength{\tabcolsep}{5pt}
\resizebox{\textwidth}{!}{%
\begin{tabular}{@{}llll@{}}
\toprule
block & reference & sites & metric \\
\midrule
lag, on site & pressure sensors along the r\'ia & Villaviciosa & lag in reaching the same level (min); \\
 & (campaign pending) & & same-instant level difference (m) \\
lag, external & gauge pair: mouth and inner (judge) & Ferrol, Westerschelde, & gauge-pair lag vs.\ lag applied (min); \\
 & & Ems-Dollard, Wadden & RMSE of the level at the judge, without \\
 & & & and with the lag (m) \\
elevation, on site & RTK-GNSS survey, development split & Villaviciosa & RMSE, slope $b$, $r$ of each map on the \\
 & & & survey (Eq.~\ref{eq:metrics}); transect profile \\
elevation, external & IGN LiDAR DTM; Vaklodingen & Spanish sites; Dutch sites & RMSE, slope $b$, $r$ (Eq.~\ref{eq:metrics}), \\
 & & & by band for the Vaklodingen \\
\bottomrule
\end{tabular}}
\end{table*}

\subsubsection{Metrics and the datum}

Lags are compared in minutes: the gauge-pair lag is the mean over
paired tidal cycles $c$ of $t^{\mathrm{in}}_c - t^{\mathrm{mouth}}_c$,
the instants at which the inner and the mouth record cross the same
level on the same limb, and it is set against the lag MAREA applies at
the inner gauge. The effect of a lag on the water level is the RMSE
between the level a product predicts there, the boundary $h_B$ read
$\tau$ minutes earlier, and the level $h^{\mathrm{in}}$ the gauge
recorded, the median of the difference (the datum) removed:
\begin{align}
\mathrm{RMSE}_h &= \sqrt{\frac{1}{N}\sum_{t}\big(d(t) - \operatorname{median}\,d\big)^{2}}, \nonumber \\
d(t) &= h_B(t - \tau) - h^{\mathrm{in}}(t),
\label{eq:judge}
\end{align}
over the $N$ ten-minute samples of 2023--2025. Elevations are compared
with Eq.~\eqref{eq:metrics}. Products are on the datum of the tide model at
the mouth, the mean sea level of EOT20 \cite{hartdavis2021}; each truth
is on its own: NAP for the Vaklodingen \cite{rws_vaklodingen}, the
ellipsoid for the RTK-GNSS survey, orthometric heights for the LiDAR. No
method can see the constant between the two, so every comparison removes
it and reports it as the bias. On the $n$ samples where product $z$ and
truth $z^{\mathrm{t}}$ are both finite, each series is centred on its own
median, $y_i = z_i - \operatorname{median}(z)$ and
$x_i = z^{\mathrm{t}}_i - \operatorname{median}(z^{\mathrm{t}})$, and
\begin{align}
\mathrm{RMSE} &= \sqrt{\frac{1}{n}\sum_{i=1}^{n}(y_i - x_i)^2}, \qquad
b = \frac{\sum_i (x_i-\bar x)(y_i-\bar y)}{\sum_i (x_i-\bar x)^2}, \nonumber \\
r &= \frac{\sum_i (x_i-\bar x)(y_i-\bar y)}
         {\sqrt{\sum_i (x_i-\bar x)^2\,\sum_i (y_i-\bar y)^2}},
\label{eq:metrics}
\end{align}
with the truth on the abscissa: the slope $b$ says whether the map
reproduces the relief (1) or compresses it (below 1), $r$ how much of it
the map explains. Only truth within the range of scene levels, widened
by $\pm 0.25$~m, is scored: no waterline method can represent a channel
bed at $-8$~m beside a flat. A $5 \times 5$ thinning of the grid, one
pixel per 100~m block, changed no ranking.

\subsubsection{The interior lag on site: pressure sensors along the r\'ia}

The R\'ia de Villaviciosa has no interior tide gauge, so its lag will be
measured directly: absolute pressure sensors at three to four stations,
submerged through the lowest spring low water, with a barometer on land,
over 29 days at 1 to 5~min. Level follows
$h = z_s + (p - p_{\mathrm{atm}})/(\rho g)$, with $\rho = 1025$~kg\,m$^{-3}$
and $z_s$ the sensor elevation from RTK-GNSS. Each station is then
scored against the mouth as the gauge pairs below. The campaign is
pending.

\subsubsection{The interior lag against gauge pairs}

Four estuaries have a gauge near the mouth and another inside, the judge
of Section~\ref{sec:sites}: Ferrol1 and Ferrol2, Vlissingen and
Terneuzen, Borkum and Delfzijl, Terschelling Noordzee and Harlingen. Two quantities need
only the two records \cite{ioc2026}, 2023--2025, 10-min grid, demeaned,
gaps over one hour left open. The first is the lag in reaching the same
level: crossings of the 25th, 50th and 75th percentile of the mouth's
level are interpolated between samples and paired with the nearest inner
crossing within three hours; their mean difference over paired cycles,
and per limb, is the gauge-pair lag. It is compared with the lag MAREA
applies at the gauge's distance and with the 95\,\% bootstrap interval of
the band containing the inner gauge, measured from band 0, which holds
the mouth gauge (\emph{no delay}: 0 everywhere). The second is the
same-instant level difference, inner minus mouth: its RMS and mean are
the error of the uniform level every published method assumes. The third
is the level predicted at the inner gauge from the mouth boundary read
$\tau$ minutes earlier, with $\tau = 0$ and with MAREA's lag there,
scored on the gauge's record as the median-centred RMSE of
Eq.~\eqref{eq:metrics}. Nothing was fitted on the inner gauge, so the
whole record is the test. The comparison method's lag, per pixel over
the flats, reaches no gauge and is scored only through its elevations.

\subsubsection{Elevation on site: the RTK-GNSS survey}

In Villaviciosa the maps are scored on the development subset of the
RTK-GNSS survey of Section~\ref{sec:sites}: 35\,\% of its 150~m blocks,
drawn once with a fixed seed and sealed before any elevation product
existed, released by a runtime guard only to the verdict script. The two
maps share record, epoch (2023--2025) and boundary (EOT20 at the mouth):
\emph{MAREA}, our inversion with the smoothed band profile at each
pixel's distance, and \emph{Granadeiro}, the method of
Section~\ref{sec:granadeiro} with its own lag map. Both are scored with
Eq.~\eqref{eq:metrics} and drawn with the survey along a straight
stretch of the walked route.

\subsubsection{Elevation external: the national LiDAR and the Vaklodingen}

At the Spanish sites, Villaviciosa and Ferrol, the reference is the 5~m
LiDAR DTM of the Instituto Geogr\'afico Nacional (PNOA),
% TODO: add a bib entry for the IGN MDT05 / PNOA-LiDAR product
resampled bilinearly onto the product grid; at the Dutch sites, the
Rijkswaterstaat Vaklodingen \cite{rws_vaklodingen}, yearly bathymetry
and intertidal topography on a 20~m grid (NAP), each pixel taking its
nearest cell's latest survey of 2019--2025. Both maps are scored with
Eq.~\eqref{eq:metrics} on the intertidal pixels each resolves and on
those both resolve, the Vaklodingen also by band of distance from the
mouth. A LiDAR flown at high water records the water surface over the
flats, so each LiDAR score comes with the water share of its truth: the
share of intertidal cells within 5~cm of the modal value.

\section{Results and discussion}\label{sec:results}

\subsection{The interior lag against the gauge pairs}\label{sec:res_lag}

\begin{table*}[t]
\centering
\caption{The interior lag against the gauge pairs, 2023--2025. Lag between
the mouth gauge and the inner gauge in reaching the same level (mean over
paired cycles and limbs, positive: the inner gauge is later); lag MAREA
applies at the inner gauge, with the 95\,\% bootstrap interval of its
band; RMS difference between the two gauge records at the same instant,
the error of copying the mouth gauge inland; and the RMSE of the level
predicted at the inner gauge from the boundary the products use (the
consensus of EOT20 and the mouth gauge), without and with the lag.
% NUMBERS: lag_pair_metrics.csv and judge_metrics.csv of each site's
% comparison notebook, run of 2026-09-23.
}
\label{tab:lag_external}
\footnotesize
\setlength{\tabcolsep}{6pt}
\begin{tabular}{@{}lrrrr@{}}
\toprule
 & lag between & lag applied & inner vs.\ mouth gauge, & level at the inner gauge from \\
site (inner gauge, km from mouth) & gauges (min) & by MAREA (min) & same instant, RMS (m) & the boundary (m): no lag $\to$ MAREA \\
\midrule
Ferrol (Ferrol2, 10.1) & $+4$ & $0$ $[0, +2]$ & 0.10 & $0.112 \to 0.112$ \\
Westerschelde (Terneuzen, 21.9) & $+19$ & $+7$ $[+4, +17]$ & 0.28 & $0.356 \to 0.311$ \\
Ems-Dollard (Delfzijl, 24.4) & $+51$ & $+63$ $[+54, +72]$ & 0.50 & $0.620 \to 0.358$ \\
Wadden, Vlie (Harlingen, 18.3) & $+88$ & $+129$ $[+123, +138]$ & 0.53 & $0.562 \to 0.303$ \\
\bottomrule
\end{tabular}
\end{table*}

% figure fig:lag_external is placed after the bibliography (main.tex)


% Numbers: lag_pair_metrics.csv and judge_metrics.csv of each site (run of
% 2026-09-23).
The gauge pairs order the sites as the archive does
(Table~\ref{tab:lag_external}, Fig.~\ref{fig:lag_external}). In Ferrol
the inner gauge lags the mouth by 4~min and one level for the whole
r\'ia errs by 0.10~m RMS; MAREA measures 0~min there, interval 0 to
$+2$, and the level error stays at 0.112~m. In the Westerschelde the
gauges are 19~min apart on both limbs; MAREA applies $+7$~min at
Terneuzen, the band's interval reaching $+17$, and the error falls from
0.356 to 0.311~m. In the Ems-Dollard the lag is 51~min, 41 rising and 61
falling; MAREA applies $+63$~min at Delfzijl, inside its interval of
$+54$ to $+72$, and the error falls from 0.620 to 0.358~m, by 42\,\%. In
the Vlie basin the gauges are 88~min apart, 91 rising and 85 falling;
MAREA applies $+129$~min at Harlingen, interval $+123$ to $+138$, 41~min
beyond the gauges, and the error still falls from 0.562 to 0.303~m, by
46\,\%. Two columns of Table~\ref{tab:lag_external} differ in kind: the
same-instant difference sets the two gauge records against each other,
whereas the level error is that of the boundary the products read, the
consensus of EOT20 and the mouth gauge, which adds the model's error to
the uniform-level assumption; hence the boundary without a lag (0.620~m
in the Ems-Dollard) is worse than the mouth gauge copied inland
(0.50~m), and MAREA with the lag (0.358~m) beats either. Where an
interior tide exists the archive measures it; where none exists it
applies the few minutes its intervals declare, at a cost of centimetres
at most.

\subsection{Elevation against the RTK survey}\label{sec:res_rtk}

\begin{table}[t]
\centering
\caption{Elevation against the RTK-GNSS survey, Villaviciosa development
subset, 2023--2025, EOT20 boundary. $n$: surveyed points the map
resolves; RMSE, slope $b$ and $r$ as in Eq.~\eqref{eq:metrics}, RTK on
the abscissa; bias: the datum offset, ellipsoidal heights against the
tide-model datum, removed before scoring.
% NUMBERS: products_marea/rtk_onsite_metrics.csv (Villaviciosa notebook,
% section 9a), run pending.
}
\label{tab:rtk}
\footnotesize
\begin{tabular}{@{}lrrrrr@{}}
\toprule
method & $n$ & RMSE (m) & $b$ & $r$ & bias (m) \\
\midrule
MAREA & 194 & 0.193 & 0.72 & 0.86 & $-52.8$ \\
Granadeiro & 163 & 0.275 & 0.39 & 0.63 & $-53.3$ \\
\bottomrule
\end{tabular}
\end{table}

% figure fig:rtk_onsite is placed after the bibliography (main.tex)



% Numbers: products_marea/rtk_onsite_metrics.csv (section 9a of the
% Villaviciosa notebook; experiments/v5_rtk_onsite.py runs the same code).
On the 194 surveyed points it resolves, MAREA reaches 0.193~m RMSE with
slope 0.72 and $r = 0.86$; the Granadeiro map reaches 0.275~m with
slope 0.39 and $r = 0.63$ on 163 points (Table~\ref{tab:rtk},
Fig.~\ref{fig:rtk_onsite}a,b). Along a straight 192-m stretch of the
walked route
(Fig.~\ref{fig:rtk_onsite}c--e) the MAREA profile follows the survey
within the pixel size, including the fall at its inland end, where the
Granadeiro profile stays flat. The lag applied on the surveyed pixels is
5~min in median and 16 at most, so Villaviciosa tests the estimator
more than the clock.

\subsection{Elevation against the external references}\label{sec:res_lidar}

\begin{figure*}[t]
\centering
\includegraphics[width=0.94\textwidth]{fig_products.png}
\caption{The products of the Wadden Sea site (Vlie inlet to Harlingen),
2023--2025, over the clearest low-water scene of the archive in grey.
(a) The MAREA elevation map on the intertidal zone, with two zoom
windows marked; (b), (d) the windows in the MAREA map (20-m grid);
(c), (e) the same windows in the Rijkswaterstaat Vaklodingen, the
official survey of the flats and channels, one value per 20-m cell taken
cell for cell without interpolation, same colour scale, datum offset
removed; (f) the hydroperiod, the fraction of time each pixel is
submerged, from map (a) and the tide series.}
\label{fig:products}
\end{figure*}

% NUMBERS: lidar_metrics.csv (Ferrol), products_marea/lidar_external_metrics.csv
% (Villaviciosa), vaklodingen_metrics.csv / vaklodingen_by_band.csv (Dutch
% sites), run of 2026-09-23; experiments/v6_lidar_external.py assembles them.
\begin{table*}[t]
\centering
\caption{Elevation against the external references: the IGN 5~m LiDAR
DTM (PNOA) at the Spanish sites, the Rijkswaterstaat Vaklodingen at the
Dutch sites. Each map on all intertidal pixels it resolves (left) and on
those both resolve (right), Eq.~\eqref{eq:metrics}, reference inside the
tide range the maps could sample. Spanish site headers give the LiDAR
diagnostic: modal value over the intertidal mask and share of cells
within 5~cm of it, the water-surface share of the truth.}
\label{tab:external}
\footnotesize
\setlength{\tabcolsep}{4pt}
\resizebox{\textwidth}{!}{%
\begin{tabular}{@{}lrrrrrrr@{}}
\toprule
 & \multicolumn{4}{c}{all resolved pixels} & \multicolumn{3}{c}{common pixels} \\
\cmidrule(lr){2-5}\cmidrule(lr){6-8}
method & $n$ & RMSE (m) & $b$ & $r$ & RMSE (m) & $b$ & $r$ \\
\midrule
\multicolumn{8}{@{}l}{\emph{R\'ia de Villaviciosa}, IGN LiDAR --- mode $+2.00$~m, 69\,\% of cells within 5~cm, above the tide range and excluded} \\
MAREA & 7\,458 & 0.803 & 0.47 & 0.33 & 0.760 & 0.42 & -- \\
Granadeiro & 4\,477 & 0.775 & 0.32 & 0.25 & 0.767 & 0.33 & -- \\
\midrule
\multicolumn{8}{@{}l}{\emph{R\'ia de Ferrol}, IGN LiDAR --- mode $0.00$~m, 78\,\% of cells within 5~cm} \\
MAREA & 51\,716 & 0.736 & 0.46 & 0.13 & 0.648 & 0.53 & -- \\
Granadeiro & 49\,262 & 0.732 & 0.36 & 0.09 & 0.644 & 0.54 & -- \\
\midrule
\multicolumn{8}{@{}l}{\emph{Westerschelde}, Vaklodingen} \\
MAREA & 62\,181 & 0.435 & 0.76 & 0.92 & 0.385 & 0.80 & 0.94 \\
Granadeiro & 45\,005 & 0.777 & 0.78 & 0.77 & 0.776 & 0.78 & 0.78 \\
\midrule
\multicolumn{8}{@{}l}{\emph{Ems-Dollard}, Vaklodingen} \\
MAREA & 412\,565 & 0.363 & 0.65 & 0.84 & 0.327 & 0.67 & 0.86 \\
Granadeiro & 98\,953 & 0.531 & 0.39 & 0.58 & 0.521 & 0.40 & 0.58 \\
\midrule
\multicolumn{8}{@{}l}{\emph{Wadden Sea (Vlie)}, Vaklodingen} \\
MAREA & 348\,894 & 0.329 & 0.81 & 0.68 & 0.230 & 1.00 & 0.83 \\
Granadeiro & 65\,576 & 0.477 & 0.65 & 0.44 & 0.427 & 0.77 & 0.53 \\
\bottomrule
\end{tabular}}
\end{table*}

% figure fig:vaklodingen_bands is placed after the bibliography (main.tex)


% Numbers: vaklodingen_metrics.csv / vaklodingen_by_band.csv (Dutch sites),
% lidar_metrics.csv (Ferrol), run of 2026-09-23. Villaviciosa pending.
Against the Vaklodingen (Table~\ref{tab:external},
Fig.~\ref{fig:vaklodingen_bands}) MAREA is the closer map at the three
Dutch sites. In the Westerschelde it scores 0.435~m RMSE with slope 0.76
and $r = 0.92$ on 62\,000 pixels, Granadeiro 0.777~m on 45\,000. In the
Ems-Dollard MAREA scores 0.363~m with slope 0.65 and $r = 0.84$ on
413\,000 pixels; Granadeiro resolves a quarter of them at 0.531~m with
slope 0.39. In the Vlie basin MAREA scores 0.329~m with slope 0.81 on
349\,000 pixels, Granadeiro 0.477~m on 66\,000; on the pixels both
resolve, 0.230 against 0.427~m. Fig.~\ref{fig:products} shows the Vlie
map beside the Vaklodingen and the hydroperiod derived from it. The lag
carries the two large sites:
inverted without it, the same pixels score 0.565~m with slope 0.24 in
the Ems-Dollard and 0.474~m in the Vlie, and beyond 28~km of the Ems the
error falls from 0.50--0.66~m to 0.28--0.42~m with the clock.
% ^ ablation of MAREA's own clock (from the same CSVs), kept as one
% sentence; not a method row in the tables.
At the Spanish sites the LiDAR is a water surface: 78\,\% of the Ferrol
cells lie within 5~cm of 0.00~m, and inside the tide range both maps
score 0.73~m there with correlations near 0.1; in Villaviciosa 69\,\% of
the 32\,120 cells lie within 5~cm of $+2.00$~m, above the tide range
and excluded, and on the 8\,630 cells left both maps score 0.78 to
0.80~m with $r$ below 0.35. The reference discriminates nothing at
either site.

\subsection{Limitations and future work}\label{sec:limitations}

\paragraph{Gain.} A wet/dry archive fixes the phase of the interior tide,
not its amplitude (Section~\ref{sec:adoption}): gain enters only as an
external prior, and a strongly amplified tide leaves it unmeasured in the
elevations.

\paragraph{Sampling.} A sun-synchronous sensor never sees $S_2$ and
samples $M_2$ at a 14.8-day alias; the lowest spring tides never enter
the archive, so the products stop at the lowest observed level.

\paragraph{Bands and smoothing.} The lag is measured per 3-km band and
applied through a kernel of the same length, so sharper phase changes
are not resolved. Lags are applied wherever fitted, noise included, at
a cost of centimetres where no lag exists.

\paragraph{Boundary.} Lags are measured against the boundary series at
the mouth, whose phase error is absorbed into the lag while the residual
level error holds the surge and harmonics it lacks; the nearest gauge
always beat the model and their consensus.

\paragraph{Epochs and truth years.} Morphology moves, so products are
fitted per epoch of two to three years; the truths have their own dates
(the Dollard, 2020).

\paragraph{The comparison method.} Granadeiro et
al.~\cite{granadeiro2021} was reimplemented from the paper with declared
substitutions; the numbers are ours, not theirs.
